\title{2020 年高考数学 全国I卷-理科}
\date{}
\maketitle

\section*{一、选择题}
本题共12小题，每小题5分，共60分。在每小题给出的四个选项中，只有一项是符合题目要求的。

\begin{question}[points=5]
若$z=1+i$，则$|z^2-2z|=$
\begin{choices}
  \item $0$
  \item $1$
  \item $\sqrt{2}$
  \item $2$
\end{choices}
\begin{solution}
\textbf{答案：}D

\textbf{详解：}由题意可得$z^2=(1+i)^2=2i$，则$z^2-2z=2i-2(1+i)=-2$，故$|z^2-2z|=|-2|=2$。
\end{solution}
\end{question}

\begin{question}[points=5]
设集合$A=\{x|x^2-4\leq 0\}$，$B=\{x|2x+a\leq 0\}$，且$A\cap B=\{x|-2\leq x\leq 1\}$，则$a=$
\begin{choices}
  \item $-4$
  \item $-2$
  \item $2$
  \item $4$
\end{choices}
\begin{solution}
\textbf{答案：}B

\textbf{详解：}由$x^2-4\leq 0$得$A=\{x|-2\leq x\leq 2\}$，由$2x+a\leq 0$得$B=\{x|x\leq -\frac{a}{2}\}$。因为$A\cap B=\{x|-2\leq x\leq 1\}$，所以$-\frac{a}{2}=1$，解得$a=-2$。
\end{solution}
\end{question}

\begin{question}[points=5]
埃及胡夫金字塔是古代世界建筑奇迹之一，它的形状可视为一个正四棱锥。以该四棱锥的高为边长的正方形面积等于该四棱锥一个侧面三角形的面积，则其侧面三角形底边上的高与底面正方形的边长的比值为
\begin{choices}
  \item $\frac{\sqrt{5}-1}{4}$
  \item $\frac{\sqrt{5}-1}{2}$
  \item $\frac{\sqrt{5}+1}{4}$
  \item $\frac{\sqrt{5}+1}{2}$
\end{choices}
\begin{solution}
\textbf{答案：}C

\textbf{详解：}如图，设底面正方形边长为$a$，侧面三角形底边上的高为$b$，则棱锥的高$h=\sqrt{b^2-\left(\frac{a}{2}\right)^2}$。由题意$h^2=\frac{1}{2}ab$，即$b^2-\frac{a^2}{4}=\frac{1}{2}ab$，化简得$4\left(\frac{b}{a}\right)^2-2\cdot\frac{b}{a}-1=0$，解得$\frac{b}{a}=\frac{1+\sqrt{5}}{4}$（负值舍去）。
\end{solution}
\end{question}

\begin{question}[points=5]
已知$A$为抛物线$C:y^2=2px(p>0)$上一点，点$A$到$C$的焦点的距离为$12$，到$y$轴的距离为$9$，则$p=$
\begin{choices}
  \item $2$
  \item $3$
  \item $6$
  \item $9$
\end{choices}
\begin{solution}
\textbf{答案：}C

\textbf{详解：}设抛物线的焦点为$F$，由抛物线的定义知$|AF|=x_A+\frac{p}{2}=12$，即$12=9+\frac{p}{2}$，解得$p=6$。
\end{solution}
\end{question}

\begin{question}[points=5]
某校一个课外学习小组为研究某作物种子的发芽率$y$和温度$x$（单位：$^\circ$C）的关系，在20个不同的温度条件下进行种子发芽实验，由实验数据$(x_i,y_i)(i=1,2,\ldots,20)$得到下面的散点图。由此散点图，在$10^\circ$C至$40^\circ$C之间，下面四个回归方程类型中最适宜作为发芽率$y$和温度$x$的回归方程类型的是

\begin{center}
  \includegraphics[width=.38\linewidth]{question-5-1.jpg}
\end{center}
\begin{choices}
  \item $y=a+bx$
  \item $y=a+bx^2$
  \item $y=a+be^x$
  \item $y=a+b\ln x$
\end{choices}
\begin{solution}
\textbf{答案：}D

\textbf{详解：}由散点图分布可知，散点图分布在一个对数函数的图象附近，因此最适合作为发芽率$y$和温度$x$的回归方程类型的是$y=a+b\ln x$。
\end{solution}
\end{question}

\begin{question}[points=5]
函数$f(x)=x^4-2x^3$的图像在点$(1,f(1))$处的切线方程为
\begin{choices}
  \item $y=-2x-1$
  \item $y=-2x+1$
  \item $y=2x-3$
  \item $y=2x+1$
\end{choices}
\begin{solution}
\textbf{答案：}B

\textbf{详解：}因为$f(x)=x^4-2x^3$，所以$f'(x)=4x^3-6x^2$，$f(1)=-1$，$f'(1)=-2$，因此所求切线方程为$y+1=-2(x-1)$，即$y=-2x+1$。
\end{solution}
\end{question}

\begin{question}[points=5]
设函数$f(x)=\cos\left(\omega x+\frac{\pi}{6}\right)$在$[-\pi,\pi]$的图像大致如下图，则$f(x)$的最小正周期为

\begin{center}
  \includegraphics[width=.38\linewidth]{question-7-1.jpg}
\end{center}
\begin{choices}
  \item $\frac{10\pi}{9}$
  \item $\frac{7\pi}{6}$
  \item $\frac{4\pi}{3}$
  \item $\frac{3\pi}{2}$
\end{choices}
\begin{solution}
\textbf{答案：}C

\textbf{详解：}由图可得函数图象过点$\left(-\frac{4\pi}{9},0\right)$，代入得$\cos\left(-\frac{4\pi}{9}\omega+\frac{\pi}{6}\right)=0$。又该点是函数图象与$x$轴负半轴的第一个交点，所以$-\frac{4\pi}{9}\omega+\frac{\pi}{6}=-\frac{\pi}{2}$，解得$\omega=\frac{3}{2}$，因此最小正周期$T=\frac{2\pi}{\omega}=\frac{4\pi}{3}$。
\end{solution}
\end{question}

\begin{question}[points=5]
$\left(x+\frac{y^2}{x}\right)(x+y)^5$的展开式中$x^3y^3$的系数为
\begin{choices}
  \item $5$
  \item $10$
  \item $15$
  \item $20$
\end{choices}
\begin{solution}
\textbf{答案：}C

\textbf{详解：}$(x+y)^5$展开式的通项$T_{r+1}=C_5^r x^{5-r}y^r$。$x\cdot T_{r+1}=C_5^r x^{6-r}y^r$，令$r=3$得$x^3y^3$系数$C_5^3=10$；$\frac{y^2}{x}\cdot T_{r+1}=C_5^r x^{4-r}y^{r+2}$，令$r=1$得$x^3y^3$系数$C_5^1=5$。故系数为$10+5=15$。
\end{solution}
\end{question}

\begin{question}[points=5]
已知$\alpha\in(0,\pi)$，且$3\cos2\alpha-8\cos\alpha=5$，则$\sin\alpha=$
\begin{choices}
  \item $\frac{\sqrt{5}}{3}$
  \item $\frac{2}{3}$
  \item $\frac{1}{3}$
  \item $\frac{\sqrt{5}}{9}$
\end{choices}
\begin{solution}
\textbf{答案：}A

\textbf{详解：}由$3\cos2\alpha-8\cos\alpha=5$得$3(2\cos^2\alpha-1)-8\cos\alpha=5$，即$6\cos^2\alpha-8\cos\alpha-8=0$，化简得$3\cos^2\alpha-4\cos\alpha-4=0$，解得$\cos\alpha=-\frac{2}{3}$或$\cos\alpha=2$（舍去）。因为$\alpha\in(0,\pi)$，所以$\sin\alpha=\sqrt{1-\cos^2\alpha}=\frac{\sqrt{5}}{3}$。
\end{solution}
\end{question}

\begin{question}[points=5]
已知$A,B,C$为球$O$的球面上的三个点，$\odot O_1$为$\triangle ABC$的外接圆，若$\odot O_1$的面积为$4\pi$，$AB=BC=AC=OO_1$，则球$O$的表面积为
\begin{choices}
  \item $64\pi$
  \item $48\pi$
  \item $36\pi$
  \item $32\pi$
\end{choices}
\begin{solution}
\textbf{答案：}A

\textbf{详解：}设$\odot O_1$半径为$r$，则$\pi r^2=4\pi$，得$r=2$。由正弦定理$AB=2r\sin60^\circ=2\sqrt{3}$，所以$OO_1=AB=2\sqrt{3}$。又$OO_1\perp$平面$ABC$，$O_1A=r=2$，所以球半径$R=OA=\sqrt{OO_1^2+r^2}=\sqrt{12+4}=4$，球表面积$S=4\pi R^2=64\pi$。
\end{solution}
\end{question}

\begin{question}[points=5]
已知$\odot M:x^2+y^2-2x-2y-2=0$，直线$l:2x+y+2=0$，$P$为$l$上的动点，过点$P$作$\odot M$的切线$PA,PB$，切点为$A,B$，当$|PM|\cdot|AB|$最小时，直线$AB$的方程为
\begin{choices}
  \item $2x-y-1=0$
  \item $2x+y-1=0$
  \item $2x-y+1=0$
  \item $2x+y+1=0$
\end{choices}
\begin{solution}
\textbf{答案：}D

\textbf{详解：}圆$M$方程化为$(x-1)^2+(y-1)^2=4$，点$M$到直线$l$的距离$d=\frac{|2\times1+1+2|}{\sqrt{2^2+1^2}}=\sqrt{5}>2$，故直线与圆相离。由圆的几何性质，四点$A,P,B,M$共圆，且$AB\perp MP$，所以$|PM|\cdot|AB|=2S_{\triangle PAM}=2\times\frac{1}{2}|PA|\cdot|AM|=2|PA|$，而$|PA|=\sqrt{|MP|^2-4}$。当$MP\perp l$时，$|MP|_{\min}=d=\sqrt{5}$，$|PA|_{\min}=1$，此时$|PM|\cdot|AB|$最小。直线$MP$的方程为$y-1=\frac{1}{2}(x-1)$，即$y=\frac{1}{2}x+\frac{1}{2}$，与$l$联立得$P(-1,0)$。以$MP$为直径的圆的方程为$(x-1)(x+1)+y(y-1)=0$，即$x^2+y^2-y-1=0$，两圆方程相减得$2x+y+1=0$，即为直线$AB$的方程。
\end{solution}
\end{question}

\begin{question}[points=5]
若$2^a+\log_2 a=4^b+2\log_4 b$，则
\begin{choices}
  \item $a>2b$
  \item $a<2b$
  \item $a>b^2$
  \item $a<b^2$
\end{choices}
\begin{solution}
\textbf{答案：}B

\textbf{详解：}设$f(x)=2^x+\log_2 x$，则$f(x)$为增函数。已知$2^a+\log_2 a=4^b+2\log_4 b=2^{2b}+\log_2 b$。计算$f(a)-f(2b)=2^a+\log_2 a-(2^{2b}+\log_2 2b)=2^{2b}+\log_2 b-(2^{2b}+\log_2 2b)=\log_2\frac{1}{2}=-1<0$，所以$f(a)<f(2b)$，故$a<2b$。
\end{solution}
\end{question}

\section*{二、填空题}
本题共4小题，每小题5分，共20分。

\begin{question}[points=5]
若$x,y$满足约束条件$\begin{cases}2x+y-2\leq 0\\ x-y-1\geq 0\\ y+1\geq 0\end{cases}$，则$z=x+7y$的最大值为。
\fillin[{$1$}]
\begin{solution}
\textbf{答案：}$1$

\textbf{详解：}绘制不等式组表示的平面区域，目标函数$z=x+7y$化为$y=-\frac{1}{7}x+\frac{1}{7}z$。直线在$y$轴上截距最大时$z$最大，联立$\begin{cases}2x+y-2=0\\ x-y-1=0\end{cases}$解得$A(1,0)$，代入得$z_{\max}=1+7\times0=1$。
\end{solution}
\end{question}

\begin{question}[points=5]
设$a,b$为单位向量，且$|a+b|=1$，则$|a-b|=$。
\fillin[{$\sqrt{3}$}]
\begin{solution}
\textbf{答案：}$\sqrt{3}$

\textbf{详解：}因为$a,b$为单位向量，所以$|a|=|b|=1$。由$|a+b|^2=a^2+2a\cdot b+b^2=2+2a\cdot b=1$，得$2a\cdot b=-1$。故$|a-b|^2=a^2-2a\cdot b+b^2=2-(-1)=3$，所以$|a-b|=\sqrt{3}$。
\end{solution}
\end{question}

\begin{question}[points=5]
已知$F$为双曲线$C:\frac{x^2}{a^2}-\frac{y^2}{b^2}=1(a>0,b>0)$的右焦点，$A$为$C$的右顶点，$B$为$C$上的点，且$BF$垂直于$x$轴。若$AB$的斜率为$3$，则$C$的离心率为。
\fillin[{$2$}]
\begin{solution}
\textbf{答案：}$2$

\textbf{详解：}由题意$AF=c-a$，$BF=\frac{b^2}{a}$。$AB$的斜率为$\frac{BF}{AF}=\frac{b^2/a}{c-a}=3$，即$b^2=3a(c-a)$。又$b^2=c^2-a^2$，代入得$c^2-a^2=3a(c-a)$，化简得$c^2-3ac+2a^2=0$，即$e^2-3e+2=0$，解得$e=2$或$e=1$（舍去）。
\end{solution}
\end{question}

\begin{question}[points=5]
如图，在三棱锥$P-ABC$的平面展开图中，$AC=1$，$AB=AD=\sqrt{3}$，$AB\perp AC$，$AB\perp AD$，$\angle CAE=30^\circ$，则$\cos\angle FCB=$。

\begin{center}
  \includegraphics[width=.38\linewidth]{question-16-1.jpg}
\end{center}
\fillin[{$-\frac{1}{4}$}]
\begin{solution}
\textbf{答案：}$-\frac{1}{4}$

\textbf{详解：}由$AB\perp AC$，$AB=\sqrt{3}$，$AC=1$，由勾股定理得$BC=2$。同理$BD=\sqrt{AB^2+AD^2}=\sqrt{6}$，所以$BF=BD=\sqrt{6}$。在$\triangle ACE$中，$AC=1$，$AE=AD=\sqrt{3}$，$\angle CAE=30^\circ$，由余弦定理得$CE^2=AC^2+AE^2-2AC\cdot AE\cos30^\circ=1+3-2\times1\times\sqrt{3}\times\frac{\sqrt{3}}{2}=1$，所以$CF=CE=1$。在$\triangle BCF$中，$BC=2$，$BF=\sqrt{6}$，$CF=1$，由余弦定理得$\cos\angle FCB=\frac{CF^2+BC^2-BF^2}{2CF\cdot BC}=\frac{1+4-6}{2\times1\times2}=-\frac{1}{4}$。
\end{solution}
\end{question}

\section*{三、解答题}
共70分。解答应写出文字说明、证明过程或演算步骤。第17~21题为必考题，每个试题考生都必须作答。第22、23题为选考题，考生根据要求作答。

\begin{question}[points=12]
设$\{a_n\}$是公比不为$1$的等比数列，$a_1$为$a_2,a_3$的等差中项。

(1)求$\{a_n\}$的公比；

(2)若$a_1=1$，求数列$\{na_n\}$的前$n$项和。
\begin{solution}
\textbf{答案：}(1) $-2$；(2) $S_n=\frac{1-(1+3n)(-2)^n}{9}$

\textbf{详解：}(1)设公比为$q$，由$2a_1=a_2+a_3$得$2a_1=a_1q+a_1q^2$，即$q^2+q-2=0$，解得$q=-2$或$q=1$（舍去）。故公比为$-2$。
(2)由$a_1=1$得$a_n=(-2)^{n-1}$，则$na_n=n(-2)^{n-1}$。
设前$n$项和为$S_n$，则$S_n=1\times1+2\times(-2)+3\times(-2)^2+\cdots+n(-2)^{n-1}$，①
$-2S_n=1\times(-2)+2\times(-2)^2+3\times(-2)^3+\cdots+(n-1)(-2)^{n-1}+n(-2)^n$，②
①-②得$3S_n=1+(-2)+(-2)^2+\cdots+(-2)^{n-1}-n(-2)^n=\frac{1-(-2)^n}{1-(-2)}-n(-2)^n=\frac{1-(-2)^n}{3}-n(-2)^n$，
所以$S_n=\frac{1-(-2)^n-3n(-2)^n}{9}=\frac{1-(1+3n)(-2)^n}{9}$。
\end{solution}
\end{question}

\begin{question}[points=12]
如图，$D$为圆锥的顶点，$O$是圆锥底面的圆心，$AE$为底面直径，$AE=AD$。$\triangle ABC$是底面的内接正三角形，$P$为$DO$上一点，$PO=\frac{\sqrt{6}}{6}DO$。

(1)证明：$PA\perp$平面$PBC$；

(2)求二面角$B-PC-E$的余弦值。

\begin{center}
  \includegraphics[width=.38\linewidth]{question-18-1.jpg}
\end{center}
\begin{solution}
\textbf{答案：}(1)证明见解析；(2)$\frac{2\sqrt{5}}{5}$

\textbf{详解：}(1)设$AE=1$，则$DO=\frac{\sqrt{3}}{2}$，$CO=BO=\frac{1}{2}$，$PO=\frac{\sqrt{6}}{6}\times\frac{\sqrt{3}}{2}=\frac{\sqrt{2}}{4}$。计算得$PC=PB=\sqrt{PO^2+OC^2}=\frac{\sqrt{6}}{4}$，$AB=\sqrt{3}OA=\frac{\sqrt{3}}{2}$。由$PA^2+PB^2=AB^2$得$PA\perp PB$，同理$PA\perp PC$，又$PC\cap PB=P$，所以$PA\perp$平面$PBC$。
(2)过$O$作$ON\parallel BC$交$AB$于点$N$，以$O$为原点，$OA$为$x$轴，$ON$为$y$轴建立空间直角坐标系。则$E(-\frac{1}{2},0,0)$，$P(0,0,\frac{\sqrt{2}}{4})$，$B(-\frac{1}{4},\frac{\sqrt{3}}{4},0)$，$C(-\frac{1}{4},-\frac{\sqrt{3}}{4},0)$。
平面$PBC$的一个法向量$\vec{n}=(\sqrt{2},0,-1)$，平面$PCE$的一个法向量$\vec{m}=(1,\frac{\sqrt{3}}{3},-\sqrt{2})$，$\cos\langle\vec{m},\vec{n}\rangle=\frac{\vec{n}\cdot\vec{m}}{|\vec{n}||\vec{m}|}=\frac{2\sqrt{2}}{\sqrt{3}\times\sqrt{\frac{10}{3}}}=\frac{2\sqrt{5}}{5}$。故二面角$B-PC-E$的余弦值为$\frac{2\sqrt{5}}{5}$。
\end{solution}
\end{question}

\begin{question}[points=12]
甲、乙、丙三位同学进行羽毛球比赛，约定赛制如下：累计负两场者被淘汰；比赛前抽签决定首先比赛的两人，另一人轮空；每场比赛的胜者与轮空者进行下一场比赛，负者下一场轮空，直至有一人被淘汰；当一人被淘汰后，剩余的两人继续比赛，直至其中一人被淘汰，另一人最终获胜，比赛结束。经抽签，甲、乙首先比赛，丙轮空。设每场比赛双方获胜的概率都为$\frac{1}{2}$。

(1)求甲连胜四场的概率；

(2)求需要进行第五场比赛的概率；

(3)求丙最终获胜的概率。
\begin{solution}
\textbf{答案：}(1)$\frac{1}{16}$；(2)$\frac{3}{4}$；(3)$\frac{7}{16}$

\textbf{详解：}(1)甲连胜四场表示甲依次战胜乙、丙、乙、丙，概率$P=\left(\frac{1}{2}\right)^4=\frac{1}{16}$。
(2)四局内结束比赛的情况有：ABAB, ACAC, BCBC, BABA（其中A表示甲输，B表示乙输，C表示丙输），每种概率$\left(\frac{1}{2}\right)^4=\frac{1}{16}$，故四局内结束的概率$P'=4\times\frac{1}{16}=\frac{1}{4}$，因此需要进行第五场比赛的概率$1-P'=\frac{3}{4}$。
(3)丙最终获胜的情况（甲赢时）包括：BCBC, ABCBC, ACBCB, BABCC, BACBC, BCACB, BCABC, BCBAC，共8种，概率$P_{\text{甲}}=\left(\frac{1}{2}\right)^4+7\times\left(\frac{1}{2}\right)^5=\frac{1}{16}+\frac{7}{32}=\frac{9}{32}$。由对称性，乙赢的概率与甲相同，故丙赢的概率$P_{\text{丙}}=1-2\times\frac{9}{32}=\frac{7}{16}$。
\end{solution}
\end{question}

\begin{question}[points=12]
已知$A,B$分别为椭圆$E:\frac{x^2}{a^2}+y^2=1(a>1)$的左、右顶点，$G$为$E$的上顶点，$\overrightarrow{AG}\cdot\overrightarrow{GB}=8$。$P$为直线$x=6$上的动点，$PA$与$E$的另一交点为$C$，$PB$与$E$的另一交点为$D$。

(1)求$E$的方程；

(2)证明：直线$CD$过定点。
\begin{solution}
\textbf{答案：}(1)$\frac{x^2}{9}+y^2=1$；(2)直线$CD$过定点$\left(\frac{3}{2},0\right)$

\textbf{详解：}(1)由椭圆方程得$A(-a,0),B(a,0),G(0,1)$，则$\overrightarrow{AG}=(a,1),\overrightarrow{GB}=(a,-1)$，$\overrightarrow{AG}\cdot\overrightarrow{GB}=a^2-1=8$，解得$a^2=9$，故椭圆方程为$\frac{x^2}{9}+y^2=1$。
(2)设$P(6,y_0)$，则直线$AP$方程为$y=\frac{y_0}{9}(x+3)$，与椭圆联立得$(y_0^2+9)x^2+6y_0^2x+9y_0^2-81=0$，由韦达定理$x_A\cdot x_C=\frac{9y_0^2-81}{y_0^2+9}$，而$x_A=-3$，得$x_C=\frac{-3y_0^2+27}{y_0^2+9}$，代入得$y_C=\frac{6y_0}{y_0^2+9}$，故$C\left(\frac{-3y_0^2+27}{y_0^2+9},\frac{6y_0}{y_0^2+9}\right)$。同理直线$PB$方程$y=\frac{y_0}{3}(x-3)$，与椭圆联立得$D\left(\frac{3y_0^2-3}{y_0^2+1},\frac{-2y_0}{y_0^2+1}\right)$。当$y_0\neq0$时，直线$CD$斜率$k=\frac{\frac{6y_0}{y_0^2+9}-\frac{-2y_0}{y_0^2+1}}{\frac{-3y_0^2+27}{y_0^2+9}-\frac{3y_0^2-3}{y_0^2+1}}=\frac{4y_0}{3(3-y_0^2)}$，直线$CD$方程为$y-\frac{-2y_0}{y_0^2+1}=\frac{4y_0}{3(3-y_0^2)}\left(x-\frac{3y_0^2-3}{y_0^2+1}\right)$，化简得$y=\frac{4y_0}{3(3-y_0^2)}\left(x-\frac{3}{2}\right)$，故过定点$\left(\frac{3}{2},0\right)$。当$y_0=0$时，$C(3,0),D(3,0)$，直线$CD$为$x=3$，也过$\left(\frac{3}{2},0\right)$。综上，直线$CD$过定点$\left(\frac{3}{2},0\right)$。
\end{solution}
\end{question}

\begin{question}[points=12]
已知函数$f(x)=e^x+ax^2-x$。

(1)当$a=1$时，讨论$f(x)$的单调性；

(2)当$x\geq 0$时，$f(x)\geq\frac{1}{2}x^3+1$，求$a$的取值范围。
\begin{solution}
\textbf{答案：}(1)当$x\in(-\infty,0)$时，$f(x)$单调递减；当$x\in(0,+\infty)$时，$f(x)$单调递增。(2)$a\in\left[\frac{7-e^2}{4},+\infty\right)$

\textbf{详解：}(1)当$a=1$时，$f(x)=e^x+x^2-x$，$f'(x)=e^x+2x-1$，$f''(x)=e^x+2>0$，故$f'(x)$单调递增。又$f'(0)=0$，所以当$x<0$时，$f'(x)<0$，$f(x)$单调递减；当$x>0$时，$f'(x)>0$，$f(x)$单调递增。
(2)由$f(x)\geq\frac{1}{2}x^3+1$得$e^x+ax^2-x\geq\frac{1}{2}x^3+1$，即$e^x-\frac{1}{2}x^3-x-1\geq -ax^2$。当$x=0$时，$1\geq1$成立。当$x>0$时，分离参数得$a\geq-\frac{e^x-\frac{1}{2}x^3-x-1}{x^2}$。令$g(x)=-\frac{e^x-\frac{1}{2}x^3-x-1}{x^2}$，则$g'(x)=-\frac{(x-2)\left(e^x-\frac{1}{2}x^2-x-1\right)}{x^3}$。令$h(x)=e^x-\frac{1}{2}x^2-x-1$，则$h'(x)=e^x-x-1$，$h''(x)=e^x-1\geq0$，故$h'(x)$单调递增，$h'(x)\geq h'(0)=0$，$h(x)$单调递增，$h(x)\geq h(0)=0$。所以当$x\in(0,2)$时，$g'(x)>0$，$g(x)$单调递增；当$x\in(2,+\infty)$时，$g'(x)<0$，$g(x)$单调递减。故$g(x)_{\max}=g(2)=\frac{7-e^2}{4}$，因此$a\geq\frac{7-e^2}{4}$。综上，$a$的取值范围是$\left[\frac{7-e^2}{4},+\infty\right)$。
\end{solution}
\end{question}

\section*{四、选考题}
请考生在第22、23题中任选一题作答。如果多做，则按所做的第一题计分。

\begin{question}[points=10]
在直角坐标系$xOy$中，曲线$C_1$的参数方程为$\begin{cases}x=\cos^k t\\ y=\sin^k t\end{cases}$（$t$为参数）。以坐标原点为极点，$x$轴正半轴为极轴建立极坐标系，曲线$C_2$的极坐标方程为$4\rho\cos\theta-16\rho\sin\theta+3=0$。

(1)当$k=1$时，$C_1$是什么曲线？

(2)当$k=4$时，求$C_1$与$C_2$的公共点的直角坐标。
\begin{solution}
\textbf{答案：}(1)曲线$C_1$表示以坐标原点为圆心，半径为1的圆；(2)$\left(\frac{1}{4},\frac{1}{4}\right)$

\textbf{详解：}(1)当$k=1$时，$C_1$参数方程为$\begin{cases}x=\cos t\\ y=\sin t\end{cases}$，消去$t$得$x^2+y^2=1$，故$C_1$表示以坐标原点为圆心，半径为1的圆。
(2)当$k=4$时，$C_1$参数方程为$\begin{cases}x=\cos^4 t\\ y=\sin^4 t\end{cases}$，$x\geq0,y\geq0$，化为$\begin{cases}\sqrt{x}=\cos^2 t\\ \sqrt{y}=\sin^2 t\end{cases}$，相加得$\sqrt{x}+\sqrt{y}=1$，即$y=x-2\sqrt{x}+1$，$0\leq x\leq1,0\leq y\leq1$。曲线$C_2$的直角坐标方程为$4x-16y+3=0$。联立$\begin{cases}\sqrt{x}+\sqrt{y}=1\\ 4x-16y+3=0\end{cases}$，消去$y$得$12x-32\sqrt{x}+13=0$，解得$\sqrt{x}=\frac{1}{2}$或$\sqrt{x}=\frac{13}{6}$（舍去），故$x=\frac{1}{4}$，$y=\frac{1}{4}$。所以公共点直角坐标为$\left(\frac{1}{4},\frac{1}{4}\right)$。
\end{solution}
\end{question}

\begin{question}[points=10]
已知函数$f(x)=|3x+1|-2|x-1|$。

(1)画出$y=f(x)$的图像；

(2)求不等式$f(x)>f(x+1)$的解集。

\begin{center}
  \includegraphics[width=.38\linewidth]{question-23-1.jpg}
\end{center}
\begin{solution}
\textbf{答案：}(1)图像见解析；(2)$\left(-\infty,-\frac{7}{6}\right)$

\textbf{详解：}(1)$f(x)=\begin{cases} x+3, & x\geq1\\ 5x-1, & -\frac{1}{3}<x<1\\ -x-3, & x\leq -\frac{1}{3}\end{cases}$，图像如图所示。
(2)将$f(x)$图像向左平移1个单位得$f(x+1)$图像。由图像可知，当$x$小于两图像交点横坐标时，$f(x)>f(x+1)$。解方程$-x-3=5(x+1)-1$得$x=-\frac{7}{6}$，故解集为$\left(-\infty,-\frac{7}{6}\right)$。
\end{solution}
\end{question}

